\scanpage{70}
\mspage{19}

une notion moins exigeante que celle de $W$-fibration forte. Si on dénote ici cette condition par b), on trouve (sous réserve que $W_\infty\subset W$), avec les notations de la p. 4

\diagram{\begin{tikzcd}[ampersand replacement=\&,column sep=small,row sep=small] \& \& \& {\vphantom b} \\ \& {a\quad W\text{-fibration ultraforte}}\arrow[d,Rightarrow] \\ \& {a'\quad W\text{-fibration forte}}\arrow[d,Rightarrow] \\ \& {b\quad W\text{-fibration}}\arrow[dl,Rightarrow]\arrow[dr,Rightarrow]\arrow[uuurr,bend right=65] \\ {a_1\quad W\text{-fibration gauche}} \& \& {a_2\quad W\text{-fibration droite}} \end{tikzcd}} \margin{Cf la bonne notion de $W$-fibration dans un contexte où on ne suppose pas $W_\infty\subset W$, dans \unsure{Prop.} et déf. 3 p. 51, 52}

\deleted{\unsure{De plus}, si} $a_1$) et $a_2$) étaient aussi appelées $\mathbb D$-fibrations cohomologiques resp. $\mathbb D$-fibrations homologiques.$^{*}$
% slanted
\margin{Non, c'est une notion encore plus forte, à \unsure{coups|coup} de car. \unsure{coexacts|exacts}, cf I p. 121.} Quand $\mathbb D$ satisfait Der 3 a$'$ (I p. 111) sur les $\mathbb D$-coefficients loc. \unsure{constants} sur les $X$ $W_\infty$-asphériques, alors les notions b, $a_1$, $a_2$ sont équivalentes. (Je présume que c'est le cas pour tous les $\mathbb D$ et tous les $W$ vraiment intéressants.) Par contre, j'ai l'impression que la notion de $W$-fibration forte, voire ultraforte, n'est guère utilisable que dans le cas $W=W_\infty$, où elle coïncide avec les précédentes. C'est dire qu'on peut probablement oublier la terminologie, et se rappeler par contre que dans le cas $W=W_\infty$, les $W$-fibrations ont des propriétés d'une force extrême, s'exprimant par les deux propriétés dégagées dans $W(4\,\mathrm a)$ et $W(4\,\mathrm b)$.

% archive sheet number: 70
